\subsection{Constants and Emergent Limits}
  \label{subsec:constants-and-emergent-limits}

  \paragraph{\texorpdfstring{$a_{\star}$}{a★} (Operational saturation threshold).}
    An effective acceleration proxy marking the onset of saturated
    response in the projected description.
    It scales inversely with the local descriptive resolution $\ell$
    according to
    \[
      a_{\star} \sim \kappa \frac{b_{\chi}}{\ell},
    \]
    where $\kappa$ is a dimensionless normalization factor.
    $a_{\star}$ is not a fundamental constant but a regime-dependent
    manifestation of the invariant transport bound $b_{\chi}$.

  \paragraph{\texorpdfstring{$b_{\chi}$}{bχ} (Invariant admissible transport bound).}
    The fundamental bound on admissible projective transport within the
    $\chi$ substrate.
    It limits the maximal rate at which relational configurations can be
    updated while preserving operational closure under projection.
    All effective saturation phenomena are manifestations of this invariant
    bound under different dimensional projections.

  \paragraph{\texorpdfstring{$c$}{c} (Effective speed of light).}
    The maximal signal propagation speed in emergent spacetime.
    $c$ is not fundamental: it is the geometric projection of the
    substrate transport bound $b_{\chi}$ through the admissible projective
    speed $c_{\chi}$ in projectable regimes.

  \paragraph{\texorpdfstring{$\hbar$}{ℏ} (Effective Planck constant).}
    An emergent quantum of action associated with projection thresholds and spectral
    granularity.
    It is not fundamental at the level of $\chi$.

  \paragraph{\texorpdfstring{$G$}{G} (Newtonian gravitational constant).}
    An emergent coupling constant arising from large-scale collective projective
    admissibility dynamics of $\chi$.
    Its value reflects structural properties rather than fundamental interaction
    strengths.

  \paragraph{\texorpdfstring{$\Lambda_{\mathrm{eff}}$}{Λeff} (Effective cosmological constant).}
    A residual large-scale effect associated with the global saturation of projective
    capacity in the late-time admissible regime.
    It reflects the asymptotic structure of the admissible ordering rather than
    a fundamental energy density of the vacuum.
